% from ega4-2-fr.tex
\oldpage[IV-2]{85}

est une extension radicielle de $K'^p$, donc si $K'$ est corps de définition
de $\overline F$, il en est de même de $K'^p$. Toutefois:
\end{remark}

\begin{proposition}[4.9.4]
\label{IV.2.4.9.4-fr}
Avec les notations de (4.9.1), supposons que $K$ soit une extension séparable
de $K'$; pour que $K'$ soit un corps de définition de $\overline T$, il faut
et il suffit que $K'$ soit un corps de définition du sous-préschéma réduit de
$X_{(K)}$ ayant $\overline T$ pour espace sous-jacent.
\end{proposition}

\begin{proof}
La condition est évidemment suffisante en vertu de (4.9.1); pour voir qu'elle
est nécessaire, notons que si $\overline T=g^{-1}(T')$, où $T'$ est une
partie fermée de $X_{(K')}$, et si $Z'$ est le sous-préschéma réduit de
$X_{(K')}$ ayant $T'$ pour espace sous-jacent, alors
$\overline Z=Z'\otimes_{K'}K=\operatorname{Spec}(K)\times_{K'}Z'$ est réduit
en vertu du critère (4.6.1, e)) et de l'hypothèse sur $K$ et a pour espace
sous-jacent $\overline T$.
\end{proof}

\begin{corollary}[4.9.5]
\label{IV.2.4.9.5-fr}
Supposons vérifiée l'une des hypothèses suivantes:
\begin{enumerate}
\item[a)] $k$ est de caractéristique $0$.
\item[b)] $K$ est une extension algébrique de $k$ et $X$ est de type fini sur
$k$.
\end{enumerate}
Alors toute partie fermée $\overline T$ de $X_{(K)}$ possède un plus petit
corps de définition, qui est une extension séparable de $k$ (et finie sur $k$
dans l'hypothèse b)).
\end{corollary}

\begin{proof}
Dans l'hypothèse a), $K$ est séparable sur chacun de ses sous-corps, et il
résulte donc de (4.9.4) que les corps de définition de $\overline T$ sont les
mêmes que ceux du sous-préschéma réduit de $X_{(K)}$ ayant $\overline T$ pour
espace sous-jacent. Le corollaire résulte alors de (4.8.11). Dans l'hypothèse
b) on sait que $K$ est une extension radicielle d'une extension séparable
$K_1$ de $k$. Alors, pour toute sous-extension $K'$ de $K$, $K'$ est
radicielle sur $K'\cap K_1$, donc, en vertu de (4.9.2), $K'$ est un corps de
définition de $\overline T$ si et seulement si $K'\cap K_1$ l'est. Cela nous
ramène au cas où $K'=K_1$, donc au cas où $K'$ est une extension algébrique
séparable de $k$, et on achève le raisonnement comme dans le cas a).
\end{proof}

\begin{proposition}[4.9.6]
\label{IV.2.4.9.6-fr}
Soient $k$ un corps, $X$ un $k$-préschéma, $K$ une extension de $k$,
$\overline U$ une partie à la fois ouverte et fermée de $X_{(K)}$,
$\overline V$ son complémentaire. Soit $K'$ une sous-extension de $K$;
alors les conditions suivantes sont équivalentes:
\begin{enumerate}
\item[a)] $K'$ est un corps de définition de la partie $\overline U$ de
$X_{(K)}$.
\item[a')] $K'$ est un corps de définition du sous-préschéma fermé de
$X_{(K)}$ induit sur l'ouvert $\overline U$.
\item[b)] $K'$ est un corps de définition de la partie $\overline V$ de
$X_{(K)}$.
\item[b')] $K'$ est un corps de définition du sous-préschéma fermé de
$X_{(K)}$ induit sur l'ouvert $\overline V$.
\end{enumerate}
\end{proposition}

\begin{proof}
Il est clair que a') entraîne a), et a) entraîne b') en vertu de (4.9.1).
Comme $\overline U$ et $\overline V$ jouent des rôles symétriques, b')
entraîne b) et b) entraîne a'), ce qui achève la démonstration.
\end{proof}

\begin{corollary}[4.9.7]
\label{IV.2.4.9.7-fr}
Sous les hypothèses de (4.9.6), il existe un plus petit corps de définition
$K'$ pour la partie $\overline U$ de $X_{(K)}$, et $K'$ est aussi le plus
petit corps de définition de la partie $\overline V$ de $X_{(K)}$ et des
sous-préschémas fermés de $X_{(K)}$ induits sur les ouverts $\overline U$ et
$\overline V$. Si en outre $X$ est de type fini sur $k$, $K'$ est une
extension finie séparable de $k$.
\end{corollary}

\begin{proof}
Compte tenu de (4.9.6) et (4.8.11), il n'y a à démontrer que la dernière
assertion. En vertu de (4.8.15, (ii)), on peut se borner au cas où $K$ est
algébriquement clos. Soit alors $\overline k$ la clôture algébrique de $k$
contenue dans $K$, et désignons par $W_i$ les composantes

